\section*{Capítulo \ref{ch:b2:mammal-reproduction} --- Reprodução sexuada dos mamíferos}

\begin{solution}{exo:b2:mammal-reproduction:1}
As espermatogônias se dividem desde a puberdade por toda a vida; as
ovogônias, só no feto. Uma \omterm{def:b2:meiosis-heredity:meiosis}{meiose} dá quatro espermatozoides, mas um
único óvulo (e três corpúsculos polares). Um espermatozoide leva cerca
de 64 dias, mais duas semanas de maturação; um óvulo fica parado na
\omterm{prop:b2:meiosis-heredity:stages}{prófase I} desde antes do nascimento até o mês de sua \omterm{def:b2:mammal-reproduction:ovary}{ovulação} ---
décadas --- e de novo na metáfase II até a fecundação. Um homem produz
$10^{8}$ espermatozoides por dia, uns $10^{12}$ ao longo da vida; uma
mulher ovula cerca de 400 óvulos.
\end{solution}

\begin{solution}{exo:b2:mammal-reproduction:2}
Ligação à \omterm{prop:b2:mammal-reproduction:fertilisation}{zona pelúcida}; \omterm{prop:b2:mammal-reproduction:fertilisation}{reação acrossômica} e digestão de um caminho;
fusão das membranas e entrada do núcleo do espermatozoide; onda de
cálcio e \omterm{prop:b2:mammal-reproduction:fertilisation}{reação cortical} (endurecimento da zona, perda de seus
receptores: o bloqueio da polispermia); término da \omterm{def:b2:meiosis-heredity:meiosis}{meiose} II e
expulsão do segundo corpúsculo polar; formação dos dois pronúcleos,
replicação do DNA, primeira mitose.
\end{solution}

\begin{solution}{exo:b2:mammal-reproduction:3}
O \omterm{def:b2:mammal-reproduction:implantation}{trofoblasto} (camada externa) forma a \omterm{def:b2:mammal-reproduction:placenta}{placenta} e os anexos; a \omterm{def:b2:mammal-reproduction:implantation}{massa
celular interna} forma o embrião. A \omterm{def:b2:mammal-reproduction:implantation}{implantação} começa no dia 6 ou 7
após a fecundação, num revestimento uterino preparado pela progesterona
do \omterm{def:b2:mammal-reproduction:ovary}{corpo lúteo}.
\end{solution}

\begin{solution}{exo:b2:mammal-reproduction:4}
Oxigênio: difusão. Glicose: transporte facilitado. Hemácias maternas:
não atravessam (os sangues nunca se misturam). IgG: transporte mediado
por receptor. Ureia: difusão, do feto para a mãe. Álcool: difunde-se
livremente. Bactérias: em geral não atravessam (algumas, como as da
sífilis e a listéria, atravessam).
\end{solution}

\begin{solution}{exo:b2:mammal-reproduction:5}
$m = 3, 1, 0.4, 0.1$: $P = 0.95, 0.63, 0.33, 0.095$. Abaixo de metade
quando $m < \ln 2 = 0.69$, isto é, abaixo de $7\times 10^{7}$
espermatozoides. A curva satura: acima de cerca de $10^{8}$,
espermatozoides a mais quase nada acrescentam; abaixo disso, a
fertilidade cai quase em proporção à contagem.
\end{solution}

\begin{solution}{exo:b2:mammal-reproduction:6}
Antes da puberdade: $7\times 10^{5}$ perdidos em 13 anos (\num{4750}
dias): cerca de 150 por dia. Depois: \num{299000} perdidos em 37 anos
(\num{13500} dias): 22 por dia; \omterm{def:b2:mammal-reproduction:ovary}{ovulação} $400/\num{13500} = 0.03$ por
dia. A atresia supera a \omterm{def:b2:mammal-reproduction:ovary}{ovulação} na razão de setecentos para um; o
estoque não se esgota pela \omterm{def:b2:mammal-reproduction:ovary}{ovulação}, mas pela morte dos ovócitos.
\end{solution}

\begin{solution}{exo:b2:mammal-reproduction:7}
Adulta: $(30/27)^{2.7} = 1.33$, $S = 1.33/2.33 = 0.57$. Fetal:
$(30/19)^{2.7} = 3.4$, $S = 3.4/4.4 = 0.77$. Às baixas pressões da
\omterm{def:b2:mammal-reproduction:placenta}{placenta}, o sangue fetal carrega um quinto a mais de oxigênio do que o
sangue adulto carregaria, e o descarrega nos tecidos a pressões em que
a hemoglobina adulta já estaria meio vazia.
\end{solution}

\begin{solution}{exo:b2:mammal-reproduction:8}
$2^{k} = 25$: $k = 4.6$ duplicações, \qty{9.3}{d}: dia 16 a 17 após a
fecundação; dia 30 a 31 após a última menstruação --- mais ou menos o
dia em que a menstruação atrasa.
\end{solution}

\begin{solution}{exo:b2:mammal-reproduction:9}
$3\times 10^{-8}\times 6\times 10^{4}\times 20/10^{-3} =
\qty{36}{mL/min}$ contra uma demanda de \qty{25}{mL/min}: a margem cai
de oito para uma e meia; o feto fica hipóxico diante de qualquer esforço
adicional, cresce devagar e pode precisar de parto antecipado.
\end{solution}

\begin{solution}{exo:b2:mammal-reproduction:10}
(a) Os \omterm{def:b2:mammal-reproduction:testis}{testículos} se formam e secretam os dois hormônios; o AMH
elimina os ductos de Müller, mas a testosterona não consegue agir: os
ductos de Wolff regridem e a genitália externa é feminina --- uma
mulher com \omterm{def:b2:mammal-reproduction:testis}{testículos} e sem útero. (b) A testosterona age, o AMH não:
um homem com epidídimo, ducto deferente e genitália masculina, mais um
útero e tubas uterinas. (c) Ovários e derivados dos ductos de Müller
(sem \omterm{def:b2:mammal-reproduction:testis}{testículo}, sem AMH); os andrógenos suprarrenais virilizam a
genitália externa em graus variados. (d) O \emph{SRY} forma
\omterm{def:b2:mammal-reproduction:testis}{testículos}, que impõem o caminho masculino: um homem com dois
cromossomos X, estéril porque faltam os genes do Y necessários à
\omterm{def:b2:mammal-reproduction:testis}{espermatogênese}.
\end{solution}

\begin{solution}{exo:b2:mammal-reproduction:11}
A mais comum (\qty{68}{\%}) é a divisão entre os dias 4 e 8: depois
de o \omterm{def:b2:mammal-reproduction:implantation}{trofoblasto} se formar (portanto uma \omterm{def:b2:mammal-reproduction:placenta}{placenta}), mas antes do âmnio
(portanto duas bolsas). Antes do dia 3, cada metade forma seu próprio
\omterm{def:b2:mammal-reproduction:implantation}{trofoblasto}, daí duas \omterm{def:b2:mammal-reproduction:placenta}{placentas} (\qty{30}{\%}); depois do dia 8, o
âmnio já está formado e é compartilhado (\qty{2}{\%}). Dividir a \omterm{def:b2:mammal-reproduction:implantation}{massa
celular interna} depois que o \omterm{def:b2:mammal-reproduction:implantation}{trofoblasto} se determinou é, pelo visto,
o evento mais fácil.
\end{solution}

\begin{solution}{exo:b2:mammal-reproduction:12}
Um marsupial investe pouco antes do nascimento: uma gestação curta, um
filhote minúsculo e uma lactação que pode ser interrompida
abandonando o filhote da bolsa se vier a seca, de modo que o risco da
mãe por tentativa é baixo e ela pode tentar de novo imediatamente. Um
placentário investe muito numa gestação longa, com um filhote grande e
bem desenvolvido que sobrevive melhor ao nascer, ao custo de uma mãe
comprometida durante meses e vulnerável nesse período. Em ambientes
produtivos e previsíveis, a maior sobrevivência dos filhotes
placentários vence; o clima errático da Austrália recompensa a opção
de abandono dos marsupiais, e seu isolamento manteve os placentários
afastados até recentemente.
\end{solution}

\begin{solution}{pb:b2:mammal-reproduction:1}
\textbf{1.} $m = 3\times 10^{8}\times 10^{-8} = 3$; $P = 1 - e^{-3} =
0.95$.
\textbf{2.} $1 - e^{-3}(1 + 3) = 0.80$: sem o bloqueio, quatro
fecundações em cinco seriam polispérmicas e o embrião triploide.
\textbf{3.} $0.15/5\times 10^{-5} = \qty{3000}{s}$, cinquenta minutos;
as contrações do útero e da tuba uterina os levam.
\textbf{4.} $300/3\times 10^{8} = 10^{-6}$: perdidos na vagina ácida,
no muco cervical, para os fagócitos do útero e na tuba errada.
\textbf{5.} De três dias antes da \omterm{def:b2:mammal-reproduction:ovary}{ovulação} a um dia depois: cerca de
quatro dias.
\textbf{6.} $m = 0.5$, $P = 0.39$; ciclos esperados $1/P = 2.5$.
\textbf{7.} As poucas mitocôndrias do espermatozoide, na peça
intermediária, são marcadas e destruídas após a entrada; o óvulo leva
cem mil.
\textbf{8.} $120/16 = 7.5$: sete divisões, $2^{7} = 128$ células ---
cerca de cem.
\textbf{9.} $\log_2 25 = 4.6$ duplicações, \qty{9.3}{d}: dia 16 a 17;
dia 30 a 31 após a menstruação.
\textbf{10.} O \omterm{def:b2:mammal-reproduction:ovary}{corpo lúteo} regride cerca de doze dias após a \omterm{def:b2:mammal-reproduction:ovary}{ovulação},
a menos que a hCG o resgate; sem ela, a progesterona cai, o
revestimento descama e o embrião se perde com a menstruação.
\textbf{11.} $\log_2 10^{5} = 16.6$ duplicações; na semana 10, a
\omterm{def:b2:mammal-reproduction:placenta}{placenta} produz sua própria progesterona e o \omterm{def:b2:mammal-reproduction:ovary}{corpo lúteo} já não é
necessário.
\textbf{12.} Os óvulos ficam na \omterm{prop:b2:meiosis-heredity:stages}{prófase I} durante décadas, e as
coesinas que mantêm os bivalentes unidos se degradam, de modo que a
não disjunção na \omterm{def:b2:meiosis-heredity:meiosis}{meiose} I aumenta muito com a idade materna; os
espermatozoides são produzidos novos em dois meses.
\textbf{13.} Depois do \omterm{def:b2:mammal-reproduction:implantation}{trofoblasto} (dia 5) e antes do âmnio (dia 8):
uma \omterm{def:b2:mammal-reproduction:placenta}{placenta}, duas bolsas.
\textbf{14.} $3\times 10^{-8}\times 1.2\times 10^{5}\times 20/3.5\times
10^{-4} = \qty{206}{mL/min}$.
\textbf{15.} $7\times 3.5 = \qty{24.5}{mL/min}$: uma margem de oito.
\textbf{16.} $500\times 0.16 = \qty{80}{mL/min}$; extração
$24.5/80 = \qty{31}{\%}$.
\textbf{17.} $5\times 3.5\times 1440 = \qty{25}{g/d}$; contidos em
\qty{28}{L} de sangue materno; fluxo diário de \qty{720}{L}: o feto
retira cerca de \qty{4}{\%} da glicose que passa.
\textbf{18.} A hemoglobina fetal está \qty{77}{\%} saturada a
\qty{30}{mmHg}, o feto tem mais hemácias e um débito cardíaco elevado,
e seus tecidos estão adaptados a trabalhar a baixa pressão --- ele
vive, na prática, no cume do Everest.
\textbf{19.} Trocas gasosas sem um pulmão funcional, nutrição e
excreção sem intestino ou rim funcionais, e os hormônios que mantêm a
gestação; ela não consegue barrar todas as toxinas --- o álcool, a
nicotina e alguns vírus atravessam.
\textbf{20.} Distensão do colo $\to$ hipotálamo $\to$ ocitocina $\to$
contração $\to$ mais distensão, até que o nascimento remova o estímulo.
Antes do termo, o útero dominado pela progesterona tem poucos
receptores de ocitocina e nenhuma junção comunicante, de modo que as
contrações não se propagam e a alça não consegue se fechar.
\textbf{21.} $750\times 2.9 = \qty{2.2}{MJ}$; custo $2.2/0.8 =
\qty{2.7}{MJ}$.
\textbf{22.} Crescimento $30\times 6 = \qty{0.18}{MJ}$, manutenção
$300\times 4 = \qty{1.2}{MJ}$: \qty{1.4}{MJ}, cerca de dois terços da
energia do leite; o restante vai para a atividade, o aquecimento e as
perdas.
\textbf{23.} $270\times 2.7 = \qty{730}{MJ}$, duas vezes e meia a
gestação.
\textbf{24.} A sucção eleva a prolactina e suprime a liberação
pulsátil do hormônio que comanda a \omterm{def:b2:mammal-reproduction:ovary}{ovulação}, de modo que a \omterm{def:b2:mammal-reproduction:ovary}{ovulação} só
volta meses depois; sem contracepção, isso espaça os nascimentos em
dois a três anos.
\textbf{25.} $P = 0.95$; teste positivo no dia 17 após a fecundação
(dia 31 do ciclo); margem de oxigênio de oito; um dia de leite custa à
mãe \qty{2.7}{MJ}.
\end{solution}
